\section[第11章\q 真正有钱人的消费模式]{第11章\\ 真正有钱人的消费模式}
\label{sec:ch11}

这天从柔伊走出33楼的电梯后，就忙翻了。这是春季特刊的截稿日，编辑部的每个人都卯足全力。大标、小标还有内文修订，全都要等她点头说了算，一个完成了下一个又来……尽管如此忙碌，柔伊的脑中依然在想着中午和芭芭拉的聚餐。她还没有回洁西卡推荐的经纪公司电话，确认要接新工作。今天她必须做这件事，而在此之前，她必须和老板谈一谈，让她知道自己快要离职了。

光是想到要对芭芭拉说这些，她都觉得难受。

下午1:00时，芭芭拉出现在柔伊的办公桌前。

\subsection{真正的有钱人，钱都花在刀口上}

外头雨停了，太阳驱散了空气中的雾气，人行道也被晒干。他们选了下城区靠近中国城的一家餐厅的户外座位。一路走过来时，她们聊了一下夏季特刊的一些构想，早在春季还没结束，夏季的构想早就开始了。

看菜单和点餐时，她们闲聊了一、两句。对这两个工作搭挡来说，这样的闲聊已经算是很长了。

简短沉默后，芭芭拉直接问：“你找我出来，是想和我说什么？”

“是的。”柔伊说，“我想找你谈谈。”话到嘴边，她又犹豫了。于是先剥开面包，沾了橄榄油。

“好的。”芭芭拉说，“要我先猜猜谈些什么吗？”

柔伊笑了：“喔，其实不用的。我就是……”她停顿了一下，放下面包，求助似的看着芭芭拉。“我有点不知该怎么说。”

芭芭拉伸出手来，放在柔伊的手臂上。“放轻松。我们共事这么久了，不用紧张的。对吗？”

柔伊点头：“有件事很怪，就是去海伦娜咖啡厅，找亨利聊的有关理财、复利之类的事，让我……”

“不安。”

“是的。”柔伊说，“非常不安。”

“为什么？”芭芭拉说，“不就是聊了聊理财的事，为什么会不安？”

“是没错。”柔伊说，“问题是聊的内容，好像是平行宇宙的事。你知道吗，有一天，我还遇见了一位石油大亨。他还说，存到人生的第一个100万有多重要。我坐在那点，一边点头一边想，喔对，人生中第一桶金，我记得当时差不多是我要拿掉牙套的时候吧！”

芭芭拉噗哧一声笑出来。

柔伊拿起沾了橄榄油的面包，却不往嘴里放。“你知道我的意思吗？我在做什么？和一位先生谈百万存款的事？我就是不应该聊着这种话题，那不是我的人生。我永远也不会变有钱的。”

芭芭拉停顿了一下，问到：“为什么？”

“因为，”柔伊尽量让自己的声音显得不那样意气用事，“因为别人是有财务自由的，芭芭拉。他们可以买昂贵的艺术品，可以在知名的经纪公司上班，可以上电视访谈。那些买我们杂志的人随时可以环游世界，想去哪就去哪，想什么时候去就什么时候去。他们是另一种人，芭芭拉，但我不是。”

“喔，”芭芭拉说。之后，又问了一句：“为什么你不是？”

“为什么我是？就是因为，我上学时贷款，贷了好大一笔，然后，我还有信用卡债，每月只能付循环利息的最低值。连吃一顿大餐对我来说，都伤筋动骨。”现在柔伊的声音已经因为有情绪而开始颤抖，“因为我在理财方面就是一个笨蛋。我没有富裕的家庭背景，这就是我啊！”

芭芭拉没有回答。

“不然，你告诉我，芭芭拉，”柔伊说，“为什么？我为什么必须变有钱？”她注意到其他桌的客人在朝这边看，但是她也不在乎。

芭芭拉平静看了看柔伊，温和地问：“可是，为什么不呢？”

柔伊叹了一口长长的气，让自己平静下来。“我也不知道为什么如此不安。”芭芭拉忍不住想笑，柔伊自己也笑了。“很神经质，对吧？”她说。

她们的餐点到了，服务生摆放餐点时，她们看着外头来来去去的人。

服务生离开后，芭芭拉说：“从我们坐下来到现在，有多少人经过我们桌子？”

柔伊估了估：“我也不知道，几百人有吧？”

芭芭拉说：“虽然我和他们并不认识，但我可以肯定，如果拦下100个人来做市调，可以发现其中大多数人没有积蓄，还有少部分人有负债。实际上，那些穿着最时尚的，可能财务状况更糟糕。他们拥有的，比他们应该要有的超过太多。”

“打肿脸充胖子。”柔伊想。

“但是我还可以告诉你一点，那就是这随机采样的100个人当中，有多少是先付钱给自己，因此存下了第一个100万的？应该有5个吧！这是我们国家的统计概率，5％，20人当中就有1个。”

“真有这回事？”柔伊问。

“是的！”芭芭拉说，“现在，考考你，看你能不能区别出哪些是那5％的人？你观察一下。”

柔伊吃了一口沙拉，继续看着来往的人。大多是职业人士，从一个会议室赶往下一个，也有游客夹在其中，从一个景点赶往下一个。

她耸了耸肩膀。“我猜不出来。我放弃。”

“你猜不出来，其实我也猜不出来。”芭芭拉说，“没有一个铁律能来区分，没有一个特征来区分出有钱人的族群。他们看起来就和其他人一样，只是他们做事的方式有一点不同。仅此而已。”

“当我说‘百万富翁’时，你想像的是什么画面？”

“一看就是有钱的样子，花大把钞票买奢侈品。”柔伊说。

芭芭拉微笑。“你这么想，也是理所当然。大多数人都是这么想的，事实却刚好相反。真正的有钱人，\textbf{大部分来该花的时候一分不少，不该花的时候一分也不多用，钱都花在刀口上}。倒是没钱的人，喜欢买花俏的东西。”

没钱的人爱买花俏的东西。柔伊想，这个她倒没听说过。

“而那些财务稳定的人，也就是百万富翁呢？”芭芭拉说。“可能就住在你家隔壁。可能是来帮你修水管的工人。”她也吃了一口午餐，“或者是咖啡厅老板。”

柔伊点头，一边继续吃沙拉。“这点没错。虽然我还是有点不敢相信，不过这是真的。我常去那家咖啡厅的老板真是个超级大富翁。”

芭芭拉又吃了一口午餐，悠悠地说：“或者，还有可能是你老板。”

柔伊从沙拉上看着芭芭拉，放下手中的刀叉，靠着椅背，看着芭芭拉好一阵，终于冒出的一句话：“什么？”

芭芭拉叹了一口气，“听着，我从不对人说这些，不要到处说，好吗？”柔伊点头，“就是亨利说的那一套，我早就开始做了。从我有收入就开始了，那时我还没认识亨利呢！”

柔伊发现自己开始结巴了。“你？怎么做到的？”

芭芭拉低头用餐。“就像我刚刚说过的，”她耸着肩膀，“你刚刚说的石油大亨，说得正确，第一桶金真的是个分水岭。首先，你会觉得自己颇有成就，也确实如此。其次是，复利让财富的累积变得更加容易。”

柔伊的头脑飞快转动，她的老板芭芭拉，一位杂志编辑总监，竟然是位百万富翁？

“可是，芭芭拉，”柔伊低声说，“那你为什么还要来杂志社上班？”

“为什么不来？我喜欢这份工作，喜欢这里的人。比方说，你啊！我们时不时一起吃饭的好时光。”芭芭拉说着，又吃了一口午餐，“和人资部的大卫谈一谈。他会协助你搞定401退休帐户的事，告诉你公司可以协助些什么。等你从本公司退休时，有一笔优渥的退休金。”

从本公司退休的时候……

这句话让柔伊回到现实，她想起了一开始找芭芭拉出来用餐的目的。她是来谈离职的，她要去洁西卡在上城区的经纪公司。

她深吸了一口气。想从芭芭拉的启发式对话中抽身，回到原本的话题上。

“其实，”她开口说，“我想谈的是……”

这时她的手机响了。柔伊看了一眼，显示的电话号码让她很惊讶。“对不起，给我1分钟，”她对芭芭拉说，接起电话。“嗨，怎么了吗？”

柔伊听了一会儿，表情僵硬地点了点头。“当然，”她说，“我今晚就到。”挂上电话后，她看着芭芭拉。

“对不起，我必须要离开了。是我妈。”
